\newpage
\phantomsection
\label{prf:relative-extremum-derivatives-equal-0-or-dont-exist}

\begin{remark*}[Return]
\hyperref[cor:necessary-condition-extremum]{Return to Theorem}
\end{remark*}

\begin{theorem*}[Necessary Condition for an Extremum]
Let $I \subseteq \mathbb{R}$ be an open interval, let $f : I \to \mathbb{R}$, and let $c \in I$. If $f$ has a relative extremum at $c$, then either $f'(c)$ does not exist, or $f'(c) = 0$.
\end{theorem*}

\begin{proof}
\textbf{Professional Standard Proof.}~\\
Let $I \subseteq \mathbb{R}$ be an open interval, let $f : I \to \mathbb{R}$, and let $c \in I$. Assume $f$ has a relative extremum at $c$. Consider two cases. If $f'(c)$ does not exist, the conclusion holds immediately. If $f'(c)$ exists, then by the Interior Extremum Theorem, $f'(c) = 0$. In either case the disjunction holds.
\end{proof}

\begin{proof}
\textbf{Detailed Learning Proof.}~\\

\textbf{Step 1.} Establish the case structure. Since $f'(c)$ either exists or does not exist, the proof considers two exhaustive cases: $f'(c)$ exists, or $f'(c)$ does not exist.

\textbf{Step 2.} Handle the case where $f'(c)$ does not exist. If $f'(c)$ does not exist, then the first disjunct of the conclusion holds immediately, making the entire disjunction true.

\textbf{Step 3.} Handle the case where $f'(c)$ exists. If $f'(c)$ exists and $f$ has a relative extremum at the interior point $c$, then by the Interior Extremum Theorem, $f'(c) = 0$. This makes the second disjunct true, so the entire disjunction holds.

\textbf{Step 4.} Conclude. Since the disjunction holds in both exhaustive cases, the theorem is proved.
\end{proof}

\begin{remark*}[Proof structure]
This is a two-case exhaustion proof based on whether the derivative exists. The first case is trivial by the definition of disjunction. The second case applies the Interior Extremum Theorem directly. The proof structure is purely logical with no calculation required.
\end{remark*}

\begin{remark*}[Dependencies]~\\
\begin{itemize}
  \item \hyperref[thm:interior-extremum-theorem]{Interior Extremum Theorem}
\end{itemize}
\end{remark*}